% =====================================================================
% MockGen (standalone package) -- long journal paper, IEEEtran journal.
%
% WRITING RULES (enforced in review):
%  * Define every acronym at first use (OData, EDMX, CDS, IR, MLP, ONNX,
%    SFT, VSIX, BAS, API, vCPU).
%  * Plain language. Short sentences. One idea per paragraph.
%  * Present tense, active voice, current-state facts only.
%  * Every measured number comes from numbers.tex -- a bare metric in
%    this file is a review defect.
%  * Every coverage number states the mode it was measured in.
%  * Label provenance (machine consensus, not human-verified) is stated
%    wherever label quality is discussed.
%  * Figure captions are self-contained and end with a takeaway.
% =====================================================================
\documentclass[journal]{IEEEtran}
\usepackage[T1]{fontenc}
\usepackage[utf8]{inputenc}
\usepackage{booktabs}
\usepackage{graphicx}
\usepackage{microtype}
\usepackage{amsmath}
\usepackage{url}
\usepackage[hidelinks]{hyperref}
\Urlmuskip=0mu plus 1mu
\emergencystretch=3em

% =====================================================================
% numbers.tex -- single source of truth for every measured number cited
% by main.tex and conference.tex in this directory.
%
% RULES
%  * One \newcommand per metric; a source-path comment above each.
%  * Prose must use these macros -- a bare number in a .tex body is a
%    review defect.
%  * A value that cannot be traced to a file in this repository or to
%    the private evidence directory must not appear here.
%  * Values recorded 2026-09-21 (standalone package 0.1.0-dev.10).
% =====================================================================

% ---------------------------------------------------------------------
% Label production. Six adjudication rounds by a three-judge panel of
% distinct models over structural field metadata only.
% source: scripts/mockserver-data-generator-training/baselines/
%         2026-09-18-v3-development-evidence.json
% ---------------------------------------------------------------------
\newcommand{\LabelFields}{23{,}741}
\newcommand{\LabelRounds}{six}
\newcommand{\LabelJudges}{three}
\newcommand{\JudgeOne}{Claude Opus 4.8}
\newcommand{\JudgeTwo}{Claude Sonnet 4.5}
\newcommand{\JudgeThree}{GPT-5.6-terra}
\newcommand{\KappaLow}{0.72}
\newcommand{\KappaHigh}{0.87}
\newcommand{\KappaFinal}{0.81}

% ---------------------------------------------------------------------
% Partitions after the family-disjoint join.
% source: ~/mockgen-private/rows/join-summary.json
% ---------------------------------------------------------------------
\newcommand{\TrainRows}{12{,}375}
\newcommand{\CalibrationRows}{4{,}991}
\newcommand{\SealedRows}{4{,}317}
\newcommand{\TrainGroups}{119}
\newcommand{\CalibrationGroups}{39}
\newcommand{\SealedServices}{41}
\newcommand{\SealedDomains}{21}
\newcommand{\TrainedClasses}{63}
\newcommand{\RoutableRoles}{44}
\newcommand{\AuxiliaryRoles}{18}

% ---------------------------------------------------------------------
% Role head architecture and calibration.
% source: ~/mockgen-private/heads/v3-head.qualified.json
% ---------------------------------------------------------------------
\newcommand{\EmbeddingDim}{384}
\newcommand{\HiddenUnits}{256}
\newcommand{\TokenBudget}{64}
\newcommand{\CalibrationAccuracy}{0.835}
\newcommand{\ConformalCoverage}{0.900}
\newcommand{\GlobalThreshold}{0.54}
\newcommand{\RoleFloor}{5}
\newcommand{\FamilyFloor}{10}
\newcommand{\PrecisionTarget}{0.95}

% ---------------------------------------------------------------------
% Sealed evaluation of the shipped role head (service-disjoint).
% source: ~/mockgen-private/heads/sealed-report.dev10.json
% ---------------------------------------------------------------------
\newcommand{\SealedDecisions}{1{,}206}
\newcommand{\SealedCorrect}{1{,}160}
\newcommand{\SealedPrecision}{96.2}
\newcommand{\SupportedCorrect}{1{,}277}
\newcommand{\SupportedTotal}{2{,}026}
\newcommand{\SupportedRecall}{63.0}
\newcommand{\StatusCorrect}{27}
\newcommand{\StatusTotal}{44}
\newcommand{\StatusRawCorrect}{33}
\newcommand{\StatusServices}{22}
\newcommand{\StatusDomains}{10}
\newcommand{\CriticalFalsePositives}{27}
\newcommand{\CriticalRate}{2.2}
\newcommand{\LexicalAccepted}{512}
\newcommand{\MetadataAccepted}{389}

% ---------------------------------------------------------------------
% Serialization study: the same frozen encoder, two renderings of the
% same field context, linear head, identical partitions.
% source: session experiment logs, scratchpad serializer-experiment.mjs
% ---------------------------------------------------------------------
\newcommand{\DelimitedAccuracy}{0.63}
\newcommand{\SentenceAccuracy}{0.84}
\newcommand{\HeadNounPrecisionBefore}{45/56}
\newcommand{\HeadNounPrecisionAfter}{91/100}

% ---------------------------------------------------------------------
% Head-capacity study: linear versus one hidden layer, same embeddings.
% source: session experiment logs, scratchpad head_capacity.py
% ---------------------------------------------------------------------
\newcommand{\LinearCalAccuracy}{0.62}
\newcommand{\MlpCalAccuracy}{0.84}
\newcommand{\LinearRoutedShare}{29}
\newcommand{\MlpRoutedShare}{74}

% ---------------------------------------------------------------------
% Field-to-value relevance verifier.
% source: ~/mockgen-private/relevance3/heads/relevance-head.qualified.json
%         ~/mockgen-private/relevance3/heads/report.json
% ---------------------------------------------------------------------
\newcommand{\RelevancePairs}{3{,}346}
\newcommand{\RelevanceTrainGroups}{19}
\newcommand{\RelevanceSealedGroups}{19}
\newcommand{\RelevancePositives}{399}
\newcommand{\RelevancePositivesAccepted}{5}
\newcommand{\RelevanceNegatives}{704}
\newcommand{\RelevanceNegativesAccepted}{6}
\newcommand{\RelevanceNegativeRate}{0.9}
\newcommand{\RelevancePositiveRate}{1.3}
\newcommand{\RelevanceLooseNegativeRate}{14}

% ---------------------------------------------------------------------
% Packaged artifacts and release identities.
% source: scripts/mockserver-data-generator-training/baselines/
%         dev10-application-modeler-1.52.0.json
%         packages/mock-data-generator/resources/models/manifest.json
% ---------------------------------------------------------------------
\newcommand{\EncoderMiB}{21.9}
\newcommand{\RoleHeadMiB}{3.4}
\newcommand{\RelevanceHeadMiB}{0.1}
\newcommand{\VocabularyMiB}{0.2}
\newcommand{\GeneratorModelMiB}{157.3}
\newcommand{\PackageMB}{128}
\newcommand{\PackageLimitMB}{200}
\newcommand{\PackageVersion}{0.1.0-dev.10}
\newcommand{\ApiVersion}{2}
\newcommand{\VsixMB}{11.4}
\newcommand{\VsixVersion}{1.52.0}

% ---------------------------------------------------------------------
% Verification state of the shipped build.
% source: session test runs, 2026-09-20
% ---------------------------------------------------------------------
\newcommand{\PackageTests}{466}
\newcommand{\TrainingTests}{80}
\newcommand{\EditorTests}{50}
\newcommand{\LintErrors}{0}
\newcommand{\ReplayServices}{four}
\newcommand{\TravelAutoMs}{2{,}845}

% ---------------------------------------------------------------------
% Corpus benchmark over every checksum-verified service of the source
% registry whose format this package parses directly. Two runs of the same
% corpus: the recognition path (tiers 0, 1 and 3, the language model given a
% 1 ms budget so it never fires) and the whole pipeline with the language
% model at its packaged service budget. Field resolution is identical in
% both; only the timings differ.
% source: scripts/mockserver-data-generator-training/baselines/
%         2026-09-21-corpus-benchmark-recognition.json  (recognition path)
%         2026-09-21-corpus-benchmark-learned.json      (whole pipeline)
% ---------------------------------------------------------------------
\newcommand{\RegistryServices}{473}
\newcommand{\CorpusServices}{131}
\newcommand{\CorpusGenerated}{131}
\newcommand{\CorpusFailed}{0}
\newcommand{\CorpusEntities}{1{,}345}
\newcommand{\CorpusFields}{11{,}529}
\newcommand{\CorpusLargestService}{897}
\newcommand{\MetadataPct}{5.96}
\newcommand{\ClassifierPct}{29.01}
\newcommand{\LexicalPct}{11.85}
\newcommand{\AbstainedPct}{53.18}
\newcommand{\ProviderBoundPct}{44.79}
\newcommand{\CorpusMedianMs}{1{,}442}
\newcommand{\CorpusPNinetyMs}{2{,}929}
\newcommand{\CorpusTotalSec}{251}
\newcommand{\CorpusFullMedianMs}{3{,}808}
\newcommand{\CorpusFullPNinetyMs}{13{,}896}
\newcommand{\CorpusFullTotalSec}{820}

% ---------------------------------------------------------------------
% Same corpus with recognition disabled (deterministic mode): the
% contribution of the learned tier, and its effect on robustness.
% source: scripts/mockserver-data-generator-training/baselines/
%         2026-09-21-corpus-benchmark-deterministic.json
% ---------------------------------------------------------------------
\newcommand{\DetGenerated}{119}
\newcommand{\DetFailed}{12}
\newcommand{\DetFields}{9{,}928}
\newcommand{\DetMetadataPct}{5.96}
\newcommand{\DetLexicalPct}{17.94}
\newcommand{\DetAbstainedPct}{76.10}
\newcommand{\DetProviderBoundPct}{22.93}
\newcommand{\DetMedianMs}{107}

% ---------------------------------------------------------------------
% Scaling benchmark: deterministic pipeline, finance service, 17 entities.
% source: scripts/mockserver-data-generator-training/baselines/
%         2026-09-19-scaling-benchmark.json
% ---------------------------------------------------------------------
\newcommand{\ScaleEntities}{17}
\newcommand{\ScaleRowsMin}{17}
\newcommand{\ScaleRowsMax}{1{,}700}
\newcommand{\ScaleMsMin}{270}
\newcommand{\ScaleMsMax}{453}
\newcommand{\ScaleMsPerRowMax}{15.9}
\newcommand{\ScaleMsPerRowMin}{0.27}

% ---------------------------------------------------------------------
% Editor availability budgets.
% source: packages/application-modeler/ide-extension/src/mockgen/runtime.ts
%         (tools-suite-standalone-mockgen)
% ---------------------------------------------------------------------
\newcommand{\AvailabilityBudgetSec}{60}
\newcommand{\GenerationBudgetSec}{300}

\newcommand{\system}{MockGen}
\newcommand{\code}[1]{\texttt{#1}}

\begin{document}

\title{\system: Realistic Offline Mock Data for Enterprise
Applications, Packaged for the Editor}

\author{Christos~Koutsiaris%
\thanks{C. Koutsiaris is with SAP P\&E -- Cloud ERP -- UX Foundation
(e-mail: christos.koutsiaris@sap.com).}}

\markboth{Preprint}{Koutsiaris: \system}

\maketitle

\begin{abstract}
Developers of enterprise applications need screens full of realistic
data long before any backend exists. Hand-written mock data is expensive
and goes stale; random generators produce values that look broken on
screen. \system{} generates complete, relationally consistent mock
datasets for an application directly from its service definition, fully
offline and deterministically. Every column is resolved by a cascade:
declared catalogs first, then a recognised semantic role, then a small
local language model, then a type-correct fallback that can never
decline. This paper describes the whole system as it ships --- the
schema representation and its adapters, the four tiers, the relational
passes that keep keys unique and references resolvable, the budget
governance that makes a small model useful on weak hardware, and the
packaging and editor integration that put it in front of developers. It
also reports the part that changed most recently: role recognition is
now a learned classifier trained on \LabelFields{} fields labelled by a
panel of \LabelJudges{} large language models that see only structural
metadata. On \CorpusServices{} real service definitions the pipeline
generates \CorpusFields{} fields with \CorpusFailed{} failures and binds
\ProviderBoundPct\% of them to a semantic provider; with recognition
disabled the same corpus binds \DetProviderBoundPct\% and
\DetFailed{} services fail outright. On \SealedServices{} services held
out from training, the classifier's own decisions are correct
\SealedPrecision\% of the time across \RoutableRoles{} roles. We also
report a component that did not work --- a field-to-value relevance
verifier --- and the shipping decision we made about it.
\end{abstract}

\begin{IEEEkeywords}
mock data, test data generation, service metadata, schema graph, OData,
offline inference, small language models, semantic role classification,
relational data generation
\end{IEEEkeywords}

\section{Introduction}
\label{sec:intro}

\IEEEPARstart{A}{frontend} developer building an enterprise application
starts from a service definition: a machine-readable contract that names
entities, columns, types, widths, and relationships. The application
must show data immediately, before any backend exists. The usual options
disappoint. Hand-written mock data is slow to produce and drifts out of
sync with the schema. Generic random generators fill a city column with
random words, so every screen looks broken. Recording real responses
needs access, permissions, and the removal of sensitive values.

Enterprise schemas make this harder. One service routinely declares
dozens of entities and hundreds of columns with exact widths and
precisions. Columns carry protocol artifacts --- draft flags, edit
shadows --- that are not business data. The schema's own annotations
encode meaning a generator should honour. Mistakes are visible: a
too-long value truncates on screen, a code that does not match its text
reads as a bug, and one empty child table makes the application look
defective.

\system{} reads the application's service definition and generates a
complete dataset for it: offline, deterministic, in seconds. Five rules
govern the design.

\begin{enumerate}
  \item \textbf{Any application, cold.} An application \system{} has
    never seen works without configuration.
  \item \textbf{No empty screens.} Lists, object pages, and child tables
    always have rows.
  \item \textbf{It must look real at a glance.} Codes match their texts,
    child rows belong to their parents, one record tells one story.
  \item \textbf{No internal failure blanks the application.} A missing
    model or a rejected value costs realism, never availability.
  \item \textbf{Determinism.} The same schema, seed, and artifacts
    produce byte-identical data.
\end{enumerate}

The architecture that delivers this is a \emph{tier cascade}
(Fig.~\ref{fig:cascade}): for each column, the first tier that can
produce a constraint-valid value wins, and the last tier cannot decline.

\begin{figure}[t]
  \centering
  \includegraphics[width=0.84\columnwidth]{figures/out/fig-cascade.pdf}
  \caption{The resolution cascade for a single column. Declared catalogs
  beat recognition, recognition beats the language model, and the typed
  fallback cannot refuse. Takeaway: the never-empty promise is
  structural, not aspirational.}
  \label{fig:cascade}
\end{figure}

\textbf{Contributions.} (i) A schema intermediate representation (IR)
and adapters that let one engine serve several service formats
(Section~\ref{sec:ir}). (ii) The four tiers and their mechanics
(Sections~\ref{sec:cascade}--\ref{sec:t3}). (iii) A learned role
recogniser trained on machine-adjudicated labels, with a routing policy
that decides per role whether the model may speak
(Sections~\ref{sec:learning}--\ref{sec:routing}). (iv) The relational
passes, determinism, and caching
(Sections~\ref{sec:relations}--\ref{sec:determinism}). (v) Packaging,
qualification, and editor integration (Section~\ref{sec:release}).
(vi) A measured characterisation of the shipped system, including a
negative result (Sections~\ref{sec:results}--\ref{sec:relevance}).

\section{Related Work}
\label{sec:related}

\textbf{Random and property-based generators} produce type-correct
values and are excellent for testing invariants, but they have no notion
of what a column means, so their output looks wrong on a screen.

\textbf{Database synthesisers} learn distributions from existing data.
They need that data; our setting has none, because the backend does not
exist yet.

\textbf{Schema matching} aligns two schemas with each other. Our problem
aligns one schema against a fixed vocabulary of generation roles. The
output is not a correspondence but a decision about which value
generator to call, and the cost of being wrong differs: a bad match is
reviewed by a human, a bad role silently writes a city into an order
number.

\textbf{Hosted language models} would classify columns well and are
unavailable at run time: the generator must work with no network, in a
locked-down development space, deterministically. We therefore use large
models offline, once, to produce labels, and ship a small model that
reproduces their decisions locally.

\section{Schema IR and Adapters}
\label{sec:ir}

Everything downstream reads one internal shape: a schema graph. Adapters
translate each supported input format into it, so the tiers never learn
protocol dialects.

The graph holds entities with their entity-set names; properties with
primitive type, width, precision, scale, key and nullability flags,
declared labels, data elements, and descriptions; relationships with
their column mappings; and annotation-derived links such as value lists,
text associations, and currency or unit companions.

Two adapters ship today. The Extensible Markup Language metadata format
(EDMX) used by Open Data Protocol (OData) services is parsed
structurally. Core Data Services (CDS) documents are parsed from their
compiled notation. The field context the recogniser reads is derived
from this graph, and it is the same object at training time and at run
time --- a property enforced by fingerprints
(Section~\ref{sec:determinism}).

\section{The Cascade at a Glance}
\label{sec:cascade}

For every column of every row, the tiers are attempted in order.

\textbf{T0, declared catalogs.} If the schema declares the allowed
values --- a value list, a code list, an enumeration --- or the
developer supplied authored rows, the value comes from there. Nothing
beats the schema's own statement of truth.

\textbf{T1, recognised roles.} Otherwise the system tries to recognise
what the column means and draws from a curated bank for that meaning.
Recognition has three sources: annotation-driven metadata rules, the
learned classifier, and audited name-based rules.

\textbf{T2, a local language model.} If no role is recognised and the
column is visible text that would otherwise read as a placeholder, a
small quantized model proposes a value under a grammar and a time
budget.

\textbf{T3, the guaranteed floor.} Otherwise a typed generator emits a
value that respects type, width, precision, and nullability. It cannot
decline.

The cascade is per column, not per table, so one entity commonly draws
from three tiers at once.

\section{T0: Authoritative Catalogs}
\label{sec:t0}

T0 is the cheapest and most trustworthy tier. A property annotated with
a value list points at a code entity; the generator reads that entity's
rows if present, and otherwise generates the code table first and draws
from it, so a code column and its text column agree by construction.

Authored data is also T0. A developer with real sample rows drops them
in, and those values win over anything the generator would produce,
while the generator still fills the columns the sample omits.

\section{T1: Recognised Roles}
\label{sec:t1}

T1 answers the question this system is really about: what does this
column mean? A role is not a free-form string. Each registered role
declares the primitive types it accepts, whether it may sit on a key,
which provider generates its values, which validator checks them, the
coherence group it belongs to, and the confidence needed to route it.
Roles are grouped into families --- status, identity, location, finance,
temporal, narrative, contact, organisation, measure --- and families
matter for both support floors and error severity. Assigning a city
where the label says postal code is an error inside the location family;
assigning a city where the label says status crosses families and counts
as critical.

\subsection{Three sources, one arbitrator}

Metadata rules fire when an annotation, data element, or explicit type
determines the role outright. Their decisions are decisive: the
classifier is not consulted.

The learned classifier handles everything else. It reads only the field
context and never sees values.

Audited name-based rules remain where a name is genuinely conclusive.
Each rule was checked against adjudicated labels; the eight that
contradicted a label with a different positive role were switched off,
while rules whose only misses were abstentions were kept, because they
carry generation value --- coherent key domains --- that labels do not
measure.

The arbitrator combines them. Metadata wins outright. If the classifier
abstains, a surviving lexical rule may fire. If the classifier routes
and a lexical rule disagrees, the classifier wins only when its
prediction set is a singleton and its confidence clears a stricter
conflict threshold; otherwise both are discarded, because two
disagreeing weak signals are worse than none. Compatibility is enforced
last: a role its registry entry forbids on keys cannot land on a key,
and a role incompatible with the column's type is rejected regardless of
source.

\subsection{Framework columns never route}

Draft-administration and action-control columns are framework
bookkeeping, not business data. They abstain unconditionally, because
inventing a plausible business value for them is worse than leaving them
to the typed floor.

\section{Learning to Recognise a Column}
\label{sec:learning}

\subsection{Labels without human annotators}

Hand-labelling enterprise columns needs domain experts and time. We
instead use a panel of \LabelJudges{} large language models from two
vendors: \JudgeOne{}, \JudgeTwo{}, and \JudgeThree{}. Each judge sees
one field's structural metadata and a frozen written guideline, and
returns a registered role or the abstention label \code{unknown}. Judges
are independent: none sees another's answer, any model prediction, or
any data value. A field is labelled when at least two of three agree;
disagreements are recorded and excluded from every partition rather than
resolved by a tie-break, which would manufacture agreement.

\LabelRounds{} rounds produced \LabelFields{} labelled fields.
Agreement, measured as Fleiss' kappa, ran from \KappaHigh{} on the first
round down to \KappaLow{} on a round aimed deliberately at the hardest
cases --- columns named \code{name}, \code{text}, \code{title} --- and
recovered to \KappaFinal{} after we clarified the guideline for exactly
those cases and re-judged them.

Every label carries its provenance: method \emph{machine panel
consensus}, the judges and their prompt fingerprints, and the flag
\code{humanVerified} set to \code{false}. We do not claim these labels
match expert human labels. We claim they are consistent, reproducible,
and honestly described.

\subsection{Partitions and leakage}

Services are grouped into families by the source registry, and a family
belongs to exactly one partition, so a sealed service shares no lineage
with anything the model learned from. Training uses \TrainRows{} rows
from \TrainGroups{} service groups; calibration uses
\CalibrationRows{} rows from \CalibrationGroups{} groups; the sealed set
is \SealedRows{} fields across \SealedServices{} services and
\SealedDomains{} domains.

Two subtler leaks are closed in the same step. Two fields can serialize
to identical text and carry different labels; those contexts are dropped
everywhere, because a context with two labels silently caps accuracy.
And a context appearing in two partitions has its lower-priority copy
dropped. Neither rule relabels anything; both only remove.

\subsection{Writing the field as language}

The classifier embeds a text rendering of the field context with a
frozen sentence encoder, then classifies the embedding
(Fig.~\ref{fig:recognition}). The rendering mattered more than anything
else we changed.

The pilot rendered a delimited record:
\code{v3|BookingStatus\_code|string|TravelBooking|...}. That is a
sensible machine format and a poor input for an encoder trained on
prose. The shipped rendering is a sentence: \code{Booking Status code
(string) in Travel Booking; data element BOOK STATUS; a code field}.
Identifiers are split into words; structural facts are stated in words;
and verbose material that consumes the \TokenBudget{}-word-piece budget
without carrying role evidence --- neighbour lists, facet objects,
annotation arrays --- is dropped from the text while remaining available
to the metadata tier. With the same encoder, labels, and head,
calibration accuracy rose from \DelimitedAccuracy{} to
\SentenceAccuracy{}. We had been preparing to fine-tune the encoder; the
measurement showed the encoder was never the bottleneck.

The final clause names the identifier's head noun. Mean pooling cannot
weigh the last word of an identifier, yet that word decides meaning:
\code{BookingStatus} is a status, \code{BookingStatusName} is a name.
Stating it raised status routing precision on calibration from
\HeadNounPrecisionBefore{} to \HeadNounPrecisionAfter{}.

\begin{figure}[t]
  \centering
  \includegraphics[width=0.82\columnwidth]{figures/out/fig-recognition.pdf}
  \caption{How one column becomes a routed role. The field context is
  rendered as a sentence, embedded by the frozen encoder, and scored by
  a small head; three gates can each send the decision back to
  abstention. Takeaway: routing is a decision about trust, not simply
  the top-scoring class.}
  \label{fig:recognition}
\end{figure}

\subsection{Head capacity}

The pilot head was a single linear layer over the embedding. With
\TrainedClasses{} classes and closely related roles, that is too little
capacity. One hidden layer of \HiddenUnits{} rectified-linear units on
the same embeddings raised calibration accuracy from
\LinearCalAccuracy{} to \MlpCalAccuracy{}, and the share of
role-bearing calibration fields the system is willing to route from
\LinearRoutedShare\% to \MlpRoutedShare\%. The shipped head is a
multilayer perceptron (MLP) of \EmbeddingDim{}$\times$\HiddenUnits{}
$\times$\TrainedClasses{}, adding \RoleHeadMiB{}\,MiB to the package and
running in microseconds per field.

\begin{table}[t]
\centering
\caption{The two design changes, measured on identical labels and the
same frozen encoder. Routed share is the fraction of role-bearing
calibration fields the system is willing to route.}
\label{tab:findings}
\footnotesize
\begin{tabular}{@{}p{0.46\columnwidth}rr@{}}
\toprule
Configuration & Cal.\ accuracy & Routed share \\
\midrule
Delimited text, linear head & \DelimitedAccuracy{} & --- \\
Sentence text, linear head & \SentenceAccuracy{} & \LinearRoutedShare\% \\
Sentence text, \HiddenUnits{}-unit head & \MlpCalAccuracy{} & \MlpRoutedShare\% \\
\bottomrule
\end{tabular}
\end{table}

\section{When the Classifier May Speak}
\label{sec:routing}

A decision routes only if it passes three gates.

\textbf{Calibration support.} The role needs at least \RoleFloor{}
correct calibration examples and its family at least \FamilyFloor{}. A
role the system cannot measure is one it must not use.

\textbf{Measured precision.} For each role we replay the runtime
acceptance rule on the calibration partition and raise that role's
threshold --- never lower it --- to the lowest confidence at which its
routed decisions reach \PrecisionTarget{} precision with enough
decisions to measure. A role that cannot reach the target at any
threshold does not route.

\textbf{Conformal set size.} Calibration yields a quantile targeting
\ConformalCoverage{} coverage; measured coverage is
\ConformalCoverage{}. A decision whose prediction set holds more than
one class abstains, whatever its confidence.

Roles failing these gates stay in the model as \emph{auxiliary} classes
with a route threshold above one: trained, never selected. This matters
more than it sounds. When we instead removed weak classes from training,
status-text columns had nowhere to go and were misrouted as status
codes. Of \TrainedClasses{} trained classes, \RoutableRoles{} route and
\AuxiliaryRoles{} are auxiliary. Coverage therefore grows with evidence
rather than with code changes: more labels promote a role from auxiliary
to routable without touching the generator.

\section{T2: A Budgeted Local Language Model}
\label{sec:t2}

When no role is recognised and the column is visible text, a small
quantized model proposes values. Three mechanisms keep it useful on two
virtual central processing units (vCPUs).

\textbf{Grammar constraint.} Decoding is constrained to a grammar
derived from the request, so output parses as rows with the expected
columns and cannot drift into prose.

\textbf{Budget governance.} Generation runs under a wall-clock budget
per entity. When the budget expires the tier yields and the cascade
continues; the run never stalls waiting for a model.

\textbf{Verification.} A separate verifier judges whether a proposed
value suits its field before it is written. Section~\ref{sec:relevance}
reports what happened when we measured it.

\section{T3: The Guaranteed Floor}
\label{sec:t3}

T3 cannot decline. It emits values respecting the declared type, width,
precision, scale, and nullability, staying unique where uniqueness is
required. Its values are plainly generic --- this is the tier that
produces \code{String 1} --- and that is the point: it turns ``no empty
screens'' from a wish into a property.

\section{Keys, References, and Caching}
\label{sec:relations}

Realism at the row level is not enough; a dataset must hold together.

\textbf{Keys.} Key columns are generated first and kept unique within
their entity, respecting width limits so a key never truncates.

\textbf{References.} A relationship's child rows draw their foreign key
values from the parent rows that were actually generated, so every
reference resolves. Child tables are populated for the parents the user
will open, so object pages are not empty.

\textbf{Coherence.} Columns in a coherence group are generated together:
a city with its country and postal code, an amount with its currency, a
quantity with its unit. A second pass recomputes derived display values
after the domains they depend on are final.

\textbf{Caching.} Generated datasets are cached under a key including
the schema, the seed, the row counts, and the fingerprints of every
model artifact. Changing any of them changes the key, so a stale dataset
is never served after a model swap.

\section{Determinism and Reproducibility}
\label{sec:determinism}

Determinism is a product promise. Generation is seeded; the classifier
performs no sampling; the text rendering is a pure function of the field
context; and the tokenizer budget is fixed, so a field that overflows it
overflows identically during training export and at run time.

Reproducibility of the models is enforced by fingerprints. The role head
records the checksum of the encoder it was trained against, the checksum
of the vocabulary, a fingerprint of the text serializer, and a
fingerprint of the role registry. At load time the runtime recomputes
all four and refuses a head whose fingerprints do not match the code it
is running inside. Changing the role vocabulary or the serialization
therefore invalidates the packaged head instead of silently degrading
it.

\section{Packaging, Qualification, and the Editor}
\label{sec:release}

\subsection{Nothing is qualified by hand}

A trained head is unqualified until a tool re-runs the full runtime
contract against it, recomputes the sealed dataset fingerprint, verifies
the label review record, checks that no developer path or unresolved
marker appears in the artifact, and only then writes the qualified
status. The packaged manifest is regenerated from file content, so a
hand-edited model file fails the release check.

\subsection{What ships where}

The editor extension carries no models (Fig.~\ref{fig:components}). It
pins the package by name, version, application programming interface
(API) version, integrity hash, and distribution hash, and resolves that
pin once into a shared module cache. The package carries the engine and
the model component: encoder \EncoderMiB{}\,MiB, role head
\RoleHeadMiB{}\,MiB, vocabulary \VocabularyMiB{}\,MiB, relevance head
\RelevanceHeadMiB{}\,MiB, and the generation model
\GeneratorModelMiB{}\,MiB. The published archive is \PackageMB{}\,MB
against a \PackageLimitMB{}\,MB registry limit. Verification is defence
in depth: the archive is checked by integrity hash on download, each
model file by checksum against the manifest, the manifest's revision is
recomputed from content, and cached executable code is re-checked before
import, with a modified file quarantined rather than used.

\begin{figure}[t]
  \centering
  \includegraphics[width=0.94\columnwidth]{figures/out/fig-components.pdf}
  \caption{What ships where. The extension holds only a pin; the package
  holds the engine and the models; a shared cache holds one verified
  copy per pinned version. Takeaway: models are distributed once and
  verified on every load, not bundled into every extension update.}
  \label{fig:components}
\end{figure}

\subsection{Availability without a stall}

One defect here is worth recording because it is easy to reproduce. The
editor decided whether to offer the action by calling the same routine
that installs the package, under a \AvailabilityBudgetSec{}-second
budget, while a first-time install of several hundred megabytes cannot
finish in that time. The menu entry therefore never appeared on a fresh
machine. The availability check now verifies platform, runtime version,
and pin, and reports a valid but uninstalled pin as available with a
pending-install flag; the install happens on first generation, inside
the \GenerationBudgetSec{}-second generation budget.

\section{Results}
\label{sec:results}

\subsection{Corpus coverage}

We run the whole pipeline over every checksum-verified service in the
source registry whose format this package parses directly:
\CorpusServices{} services of \RegistryServices{} registered, covering
\CorpusEntities{} entities and \CorpusFields{} fields, at two rows per
entity. All \CorpusGenerated{} generate with \CorpusFailed{} failures,
including a service with \CorpusLargestService{} fields.

Of those fields, \ProviderBoundPct\% bind to a semantic provider: the
classifier contributes \ClassifierPct\%, metadata rules \MetadataPct\%,
and audited lexical rules \LexicalPct\%. The remaining \AbstainedPct\%
abstain and are served by the lower tiers.

Running the same corpus with recognition disabled isolates what the
learned tier adds (Fig.~\ref{fig:resolution}). Semantic binding falls
from \ProviderBoundPct\% to \DetProviderBoundPct\%, and --- the result
we did not anticipate --- \DetFailed{} of \CorpusServices{} services
fail outright, because without a recognised role some code columns
cannot satisfy their own semantic validation. Recognition is therefore
not only a realism feature but a robustness one.

\begin{figure}[t]
  \centering
  \includegraphics[width=\columnwidth]{figures/out/fig-resolution-bars.pdf}
  \caption{How each field of the \CorpusServices{}-service corpus is
  resolved, with and without the learned recognition tier. Takeaway: the
  classifier roughly doubles the share of fields bound to a semantic
  provider, from \DetProviderBoundPct\% to \ProviderBoundPct\%.}
  \label{fig:resolution}
\end{figure}

\begin{table}[t]
\centering
\caption{Corpus profile over \CorpusServices{} service definitions, two
rows per entity. Percentages are shares of all fields in that mode. Both
runs isolate recognition, so the language model is given a budget it
cannot use; the times are therefore the recognition path, not the whole
pipeline.}
\label{tab:corpus}
\footnotesize
\begin{tabular}{@{}p{0.46\columnwidth}rr@{}}
\toprule
Measure & Recognition on & off \\
\midrule
Services generating & \CorpusGenerated{} & \DetGenerated{} \\
Failures & \CorpusFailed{} & \DetFailed{} \\
Fields & \CorpusFields{} & \DetFields{} \\
\quad bound to a provider & \ProviderBoundPct\% & \DetProviderBoundPct\% \\
\quad \quad classifier & \ClassifierPct\% & --- \\
\quad \quad metadata rules & \MetadataPct\% & \DetMetadataPct\% \\
\quad \quad lexical rules & \LexicalPct\% & \DetLexicalPct\% \\
\quad abstained & \AbstainedPct\% & \DetAbstainedPct\% \\
Median service time & \CorpusMedianMs{}\,ms & \DetMedianMs{}\,ms \\
\bottomrule
\end{tabular}
\end{table}

\subsection{Recognition quality on unseen services}

Corpus shares say how often the classifier speaks; they do not say
whether it is right. For that we use the sealed set, which is
family-disjoint from training and calibration and is scored through the
real runtime path rather than a notebook metric.

On \SealedServices{} services and \SealedDomains{} domains the model
never saw, it makes \SealedDecisions{} routing decisions of which
\SealedCorrect{} are correct: \SealedPrecision\% precision across
\RoutableRoles{} roles. Recall on supported fields --- those whose
expected role is one the head routes and which carry no decisive
annotation --- is \SupportedCorrect{}/\SupportedTotal{}
(\SupportedRecall\%). Status detection, the hardest single case because
status columns rarely carry annotations, is
\StatusCorrect{}/\StatusTotal{} across \StatusServices{} services and
\StatusDomains{} domains. The critical false-positive rate --- a role
from a different family, or a role on a field reviewers marked unknown
--- is \CriticalRate\% of accepted decisions.

\begin{table}[t]
\centering
\caption{Sealed evaluation on \SealedServices{} services never used for
training or calibration (\SealedRows{} fields). Precision counts the
classifier's own accepted decisions; deterministic metadata and lexical
decisions are listed separately because they are rules, not model
output. Takeaway: when the model speaks it is right
\SealedPrecision\% of the time.}
\label{tab:sealed}
\footnotesize
\begin{tabular}{@{}p{0.56\columnwidth}r@{}}
\toprule
Measure & Value \\
\midrule
Classifier routing decisions & \SealedDecisions{} \\
\quad correct & \SealedCorrect{} (\SealedPrecision\%) \\
\quad critical false positives & \CriticalFalsePositives{} (\CriticalRate\%) \\
Distinct roles routed & \RoutableRoles{} \\
Supported-field recall & \SupportedRecall\% \\
Status fields detected & \StatusCorrect{}/\StatusTotal{} \\
\quad classifier's top label correct & \StatusRawCorrect{}/\StatusTotal{} \\
Decisions from metadata rules & \MetadataAccepted{} \\
Decisions from lexical rules & \LexicalAccepted{} \\
\bottomrule
\end{tabular}
\end{table}

Two numbers in Table~\ref{tab:sealed} belong together. Status recall is
\StatusCorrect{}/\StatusTotal{}, while the classifier's top label is
correct on \StatusRawCorrect{}/\StatusTotal{}. The first gap is the
routing policy declining what it cannot justify; the second is the model
being wrong. Only the second closes with more labels.

The residual errors concentrate on one boundary: a generic caption
against a specific named thing. The model answers \code{description}
where the panel said \code{currency\_name}, \code{country\_name}, or
\code{org\_name} --- the same boundary on which the judges themselves
needed a guideline revision.

\subsection{Speed}

Generation time is dominated by fixed costs, not dataset size
(Fig.~\ref{fig:scaling}). On a \ScaleEntities{}-entity service with
recognition disabled, \ScaleRowsMin{} rows cost \ScaleMsMin{}\,ms and
\ScaleRowsMax{} rows cost \ScaleMsMax{}\,ms --- from
\ScaleMsPerRowMax{}\,ms per row down to \ScaleMsPerRowMin{}\,ms per row.
Across the corpus the median service completes in \DetMedianMs{}\,ms
without recognition and \CorpusMedianMs{}\,ms with it, the difference
being model load and per-field inference; the 90th percentile with
recognition is \CorpusPNinetyMs{}\,ms and the whole corpus takes
\CorpusTotalSec{}\,s. Those two runs isolate recognition, so both give
the language model a budget it cannot use. Letting it run at its
packaged budget raises the median to \CorpusFullMedianMs{}\,ms and the
corpus to \CorpusFullTotalSec{}\,s, with the same fields resolved by
the same tiers --- the extra time buys display text for columns no
catalog and no role could fill. The reference application, with every
tier enabled, generates in \TravelAutoMs{}\,ms.

\begin{figure}[t]
  \centering
  \includegraphics[width=\columnwidth]{figures/out/fig-scaling.pdf}
  \caption{Deterministic-pipeline scaling on a \ScaleEntities{}-entity
  service: total time rises from \ScaleMsMin{}\,ms to
  \ScaleMsMax{}\,ms while rows rise from \ScaleRowsMin{} to
  \ScaleRowsMax{}, so cost per row falls by roughly sixty times.
  Takeaway: everything except the language model is effectively free at
  realistic dataset sizes.}
  \label{fig:scaling}
\end{figure}

\subsection{End-to-end behaviour}

A replay across \ReplayServices{} services covering both protocols
completes in every mode the product supports. One case fails by design:
deterministic mode refuses a service whose display domains are
synthetic, because no verifier is available in that mode to check
proposed captions.

The shipped build passes \PackageTests{} package tests,
\TrainingTests{} training-tool tests, and \EditorTests{} editor tests
with \LintErrors{} lint errors, and produces an extension of
\VsixMB{}\,MB (version \VsixVersion{}) pinned to package
\PackageVersion{} at API version \ApiVersion{}.

\section{A Component That Did Not Work}
\label{sec:relevance}

\subsection{What it was for}

Before a language-model-proposed display value is written, a verifier
should judge whether it suits the field. The intended gate was that the
verifier accepts at least 80\% of valid candidates while accepting at
most 1\% of deliberately wrong ones.

\subsection{What we measured}

We built the corpus: \RelevancePairs{} field-value pairs judged by the
same panel, drawn from \RelevanceTrainGroups{} approved value sources
with \RelevanceSealedGroups{} further sources held out. We tried the
frozen encoder with a logistic head, a fine-tuned bi-encoder, and a
cross-encoder. None separated the classes. At the 1\% wrong-value budget
the verifier accepts \RelevancePositiveRate\% of valid candidates;
loosening the budget enough to be useful admits about
\RelevanceLooseNegativeRate\% of clearly-wrong values.

Expanding the corpus from seven service groups to
\RelevanceTrainGroups{} made the task harder rather than easier, and the
panel judged a sixth of our supposed ``wrong'' values as acceptable.
Both point the same way: judging whether a plausible value belongs to an
unseen service is not learnable from these features, and the task as
posed was not well defined.

\subsection{What we ship}

We ship the strict operating point:
\RelevancePositivesAccepted{}/\RelevancePositives{} valid candidates
accepted and \RelevanceNegativesAccepted{}/\RelevanceNegatives{} wrong
ones (\RelevanceNegativeRate\%). In practice it rejects nearly
everything proposed.

That made the generator refuse to produce data at all for applications
with synthetic display domains, because the pipeline treated
``unverified'' as fatal. We changed it: when the verifier declines every
candidate, generation keeps the deterministic values, marks the resource
as unverified, and emits a diagnostic. The developer gets a complete
dataset with honest labelling instead of an error. The guardrail for a
\emph{missing} verifier is unchanged, so deterministic mode still
refuses synthetic domains by design.

We report this as a negative result rather than tuning a number. A
verifier that accepts almost nothing is safe and nearly useless, and
saying so is more useful to the next implementer.

\section{Discussion}
\label{sec:discussion}

\subsection{Three mistakes, one lesson}

The pilot recogniser routed nothing useful, and none of the three causes
was the model. The input text was written for a machine rather than for
the encoder. The head had too little capacity for the class count. The
routing threshold came from a registry default that ignored the head's
own calibration, so calibrated decisions between 0.6 and 0.9 never
routed.

Each was invisible in aggregate accuracy and obvious once every sealed
decision recorded its confidence, its threshold, its prediction-set
size, and its abstention reason. Instrumentation turned one mystery into
three separate, fixable defects.

\subsection{Machine labels buy scale, not truth}

The panel produced \LabelFields{} labels for the cost of compute, and
that is why the system routes \RoutableRoles{} roles rather than a
handful. But the panel is also the ceiling: where judges disagree, the
model cannot do better, and our residual errors sit exactly where
inter-judge agreement was lowest.

\subsection{Abstention as a feature}

Because a lower tier always produces a valid value, abstention costs
realism but never availability. That asymmetry justifies gates that
would be too conservative in a conventional classifier, and it is why we
report precision on the classifier's own decisions: the product question
is not ``how often is the model right?'' but ``when it speaks, can we
trust it?''

\section{Limitations and Threats to Validity}
\label{sec:limitations}

\textbf{Label bias.} Judges are language models; a blind spot shared by
all three passes the majority vote unnoticed. We reduce but do not
remove this by using two vendors and by recording the share of decisions
resting on the column name alone.

\textbf{Coverage measurement.} Corpus numbers count fields bound to a
semantic provider, not fields a human would judge realistic. They are a
structural measure; the sealed evaluation is the quality measure.

\textbf{Corpus scope.} The corpus is \CorpusServices{} services parsable
in the two supported formats, out of \RegistryServices{} registered;
services stored only as derived graphs are excluded. The sealed set is
family-disjoint but drawn from the same registry, so it measures
generalisation across services rather than across ecosystems.

\textbf{Recall.} The head routes \SupportedRecall\% of supported sealed
fields; many columns still fall through to lower tiers.

\textbf{The verifier.} It is conservative to the point of near
inactivity, and the relevance task as posed was not well defined.

\textbf{Gate movement.} Thresholds were revised once, with the package
owner's approval, when the sealed set grew from 45 fields to
\SealedRows{}. Revising a gate after seeing results risks fitting the
gate to the model; both old and new values, the reason, and the approval
are recorded in the release baseline so a reader can apply the original
gate.

\section{Future Work}
\label{sec:future}

Three directions follow. More labelled examples on the
name-versus-description boundary, where both panel and model are
weakest. A verifier defined differently --- checking value shape against
a field's own observed domain rather than plausibility in the abstract.
And a calibration split chosen so code-list services appear in the
calibration partition, which would let several near-miss roles cross the
support floor.

\section{Conclusion}
\label{sec:conclusion}

\system{} turns a service definition into a complete, coherent,
deterministic mock dataset without a backend, a network, or a human
labelling effort. The cascade makes empty screens structurally
impossible; the recogniser makes the values worth looking at. On
\CorpusServices{} real services it generates \CorpusFields{} fields with
\CorpusFailed{} failures and binds \ProviderBoundPct\% of them to a
semantic provider, against \DetProviderBoundPct\% and \DetFailed{}
failures without recognition; on \SealedServices{} unseen services its
learned decisions are correct \SealedPrecision\% of the time across
\RoutableRoles{} roles. The engineering lessons are transferable and
unglamorous: write a model's input as language, size the head to the
class count, let calibration set the thresholds, and instrument the
evaluation so a failure tells you which of those is wrong.

\section*{Glossary}

\begin{description}
  \item[API] Application programming interface.
  \item[BAS] SAP Business Application Studio, a browser-based
    development environment.
  \item[CDS] Core Data Services, SAP's schema definition language.
  \item[Conformal prediction] A calibration method producing a set of
    plausible classes at a target coverage level.
  \item[EDMX] The XML document format for OData service metadata.
  \item[Fleiss' kappa] An agreement measure for more than two raters,
    corrected for chance.
  \item[IR] Intermediate representation; here, the schema graph.
  \item[MLP] Multilayer perceptron, a network with a hidden layer.
  \item[OData] Open Data Protocol, used by SAP Fiori applications.
  \item[ONNX] Open Neural Network Exchange, the packaged model format.
  \item[SFT] Supervised fine-tuning; here, the local language model
    tier.
  \item[vCPU] Virtual central processing unit.
  \item[VSIX] The package format for a code editor extension.
\end{description}

\end{document}
